% Page budget: Section 4, together with the next source fragment, 6--7 pages.
\section{Optimal Innovation-Safe Compression}
\label{sec:compression}

Starting from a verified source library \(L\), compression asks how little
resource burden can be retained without changing what the controller can
operate now or help generate later.  Section~3 showed that these productive
requirements are summarized by \(K_L=(F_L,C_L)\).  Throughout this section,
\emph{safe} means source-relative preservation of that exact pair,
\emph{global minimum} means minimum burden over every safe sublibrary of the
source, and an \emph{optimizer} in the capacity or penalized problems means a
maximizer over the full outer-certified family \(\mathfrak L_\theta\).  The
sequence moves from exact representation and certified deletion to the
local--global boundary, the frontier-only loss, and the two outer selection
problems.

\subsection{Minimum-resource safe representation}

\emph{Economic question.}  Among all retained sublibraries that preserve both
current operating capability and future generative capability, what is the
minimum retention burden and which representations attain it?

Recall the source-relative family
\[
 \mathcal P_{\mathrm{safe}}(L)
 =\{L'\subseteq L:s_0\in L',\ K_{L'}=K_L\}.
\]
It is a projection fiber restricted to sublibraries of \(L\).  Its global
minimum-weight correspondence is
\begin{equation}\label{eq:global-safe-optimizer}
 \mathcal M_{\mathrm{safe}}(L)
 :=\operatorname*{arg\,min}_{L'\in\mathcal P_{\mathrm{safe}}(L)}W(L').
\end{equation}

\begin{theorem}[Minimum-resource innovation-safe representation]
\label{thm:minimum-safe-compression}
For every finite source library, \(\mathcal P_{\mathrm{safe}}(L)\) is
nonempty and finite, so \(\mathcal M_{\mathrm{safe}}(L)\ne\varnothing\).
Every two safe-feasible libraries have the same compressed state and hence
the same finite-horizon productive values.  Under the discounted contraction,
they also have the same stationary productive value and the same complete set
of optimal stationary actions.
\end{theorem}

The theorem proves existence and productive preservation.  The source itself
makes the feasible family nonempty, finiteness makes the minimum attained, and
\(K_{L'}=K_L\) carries every feasible representation to the same productive
process.  Contraction is needed only for stationary value and action
preservation; detectability and probabilistic independence are not used.
This is not a uniqueness or pruning theorem: it neither singles out one
minimum-weight representation nor says that a deletion rule reaches one.
Appendix~B gives the finite argument.

For identity closure, exact safe feasibility has a useful binary cover form.
Let
\(x_s\in\{0,1\}\) indicate retention, fix \(x_{s_0}=1\), and define the source
attainers and module carriers by
\[
 A_b:=\{s\in L:j_s(b)=F_L(b)\},\qquad
 R_m:=\{s\in L:m\in\mods(s)\}.
\]
Safe feasibility is then equivalent to
\begin{equation}\label{eq:identity-safe-cover}
 \sum_{s\in A_b}x_s\ge1\quad(b\in B),\qquad
 \sum_{s\in R_m}x_s\ge1\quad(m\in C_L),
\end{equation}
with objective \(\min\sum_sw_sx_s\).  For general closure, the second family
of rows is replaced by the exact equality
\(\cl(\bigcup_{s:x_s=1}\mods(s))=C_L\); independent carrier rows would miss
closure complementarities.

\subsection{Innovation-safe deletion certificate}

\emph{Economic question.}  When can one retained strategy be removed without
changing either current operating opportunities or future research
opportunities?

\begin{definition}[Deletion redundancy]
\label{def:deletion-redundancy}
For \(s\in L\setminus\{s_0\}\), write \(L^{-s}=L\setminus\{s\}\).  Deletion
has two distinct certificates:
\begin{align}
 F_{L^{-s}}=F_L
   &\quad\text{(operational redundancy)},
   \label{eq:operational-redundancy}\\
 C_{L^{-s}}=C_L
   &\quad\text{(generative redundancy)}.
   \label{eq:generative-redundancy}
\end{align}
The first compares the entire belief-indexed frontier, not only the current
payoff.  The second compares closed capabilities, not only duplicated raw
module rows.
\end{definition}

To state the observable converse, let \(L\sim_{\DI}L'\) mean that the two
libraries have the same frontier and, for every belief and project, the same
availability-tagged cost, duration, joint terminal-belief/next-compressed-state
law, and incumbent operating-reward block.
\begin{definition}[Dynamic innovation equivalence]
\label{def:dynamic-innovation-equivalence}
With \(\langle z\rangle_{A(K)}=\operatorname{some}(z)\) when \(q\in A(K)\)
and \(\operatorname{none}\) otherwise, \(L\sim_{\DI}L'\) means equality of
the five observations
\begin{equation}\label{eq:di-definition}
 \left(F_L(b),
 \langle\kappa_q(b,K_L)\rangle_{A(K_L)},
 \langle d_q\rangle_{A(K_L)},
 \langle\overline{\mathcal Q}_q(b,K_L)\rangle_{A(K_L)},
 \langle R_q^{\op}(b,K_L)\rangle_{A(K_L)}\right)
\end{equation}
with their \(L'\) counterparts for every \(b,q\).
\end{definition}

\begin{assumption}[Observable raw closure detectability]
\label{asm:raw-closure-detectability}
If two realizable states have the same frontier but different closures, some
project changes an availability-tagged cost, duration, or joint terminal law
after generation, admission, raw update, and projection.
\end{assumption}

Detectability is used only for necessity.  Without it, a module may be
present but behaviorally silent in every available experiment.

\begin{theorem}[Frontier--closure characterization and safe deletion]
\label{thm:raw-frontier-closure-characterization}
\label{thm:unified-safe-deletion}
Under the raw factorization of Section~3,
\[
 K_L=K_{L'}\quad\Longrightarrow\quad L\sim_{\DI}L'.
\]
Under \cref{asm:raw-closure-detectability}, the implication is an
equivalence.  Consequently, for every \(s\in L\setminus\{s_0\}\),
\begin{equation}\label{eq:t3-state}
 K_{L^{-s}}=K_L
 \quad\Longleftrightarrow\quad
 \left\{
 \begin{array}{l}
 F_{L^{-s}}=F_L,\\[-2pt]
 C_{L^{-s}}=C_L.
 \end{array}\right.
\end{equation}
The two equalities certify innovation-safe deletion and preserve every
finite-horizon productive value.  Under detectability they are also necessary
for dynamic innovation equivalence; under the discounted contraction they
preserve stationary value, all stationary action comparisons, and the full
optimal-action correspondence.
\end{theorem}

The theorem proves that the two deletion tests are exactly compressed-state
preservation.  Under the Section~3 factorization, they are sufficient for
dynamic innovation equivalence and productive-value preservation; under
detectability, they are also necessary for the observable equivalence.  The
certificate is local to one candidate in the current library.  If generation
could inspect hidden raw identifiers, sufficiency would fail; without
detectability, only the necessity direction would fail.  Neither direction
factors the joint completion law, and contraction enters only after state
equality to obtain the stationary clauses.  The bridge in
\cref{thm:normalized-frontier-pruning-loss} shows where the operational-only
test fails: operational redundancy alone does not imply generative redundancy.
Appendix~A records the assumption-boundary examples and quotient refinement.

\subsection{Rechecked pruning}
\label{subsec:safe-deletion}

\emph{Economic question.}  How can repeated deletions remain certified when
removing one strategy may make another strategy essential?

The deletion certificate yields the following reduction routine.
\begin{center}
\fbox{\begin{minipage}{0.91\linewidth}
\textbf{Rechecked innovation-safe pruning.}
Start at \(L_0=L\).  At step \(i\), choose an active \(s\in L_i\), recompute
\(F_{L_i^{-s}}\) and \(C_{L_i^{-s}}\), and erase \(s\) only if both equal
their current values.  Restart the scan after every accepted deletion and
stop only after no retained active strategy passes both tests.  Return the
endpoint and the complete deletion trace.
\end{minipage}}
\end{center}

\begin{corollary}[Rechecked pruning certificate]
\label{cor:stepwise-safe-compression}
For every rechecked trace
\(L_0\supseteq L_1\supseteq\cdots\supseteq L_m\),
\[
 K_{L_m}=K_{L_0},\qquad L_m\sim_{\DI}L_0.
\]
Every accepted deletion is safe, the endpoint remains in
\(\mathcal P_{\mathrm{safe}}(L_0)\), and all productive values are preserved.
If the final scan is complete, its failure to find another safe erasure is a
certificate of one-deletion irreducibility.
\end{corollary}

The corollary proves feasibility of the entire path by composing the
rechecked compressed-state equalities.  A valid trace alone does not prove
that the endpoint is irreducible; that conclusion additionally requires the
complete final scan specified above.

Call a safe library \(E\) \emph{one-deletion irreducible} when no retained
active \(s\) satisfies \(K_{E^{-s}}=K_L\).  Call \(E\)
\emph{inclusion-wise irreducible} when no proper sublibrary of \(E\) belongs
to \(\mathcal P_{\mathrm{safe}}(L)\).  In this monotone exact-safe fiber the
two notions coincide: any proper safe sublibrary implies a safe one-element
deletion by frontier and closure monotonicity.  A complete rechecked scan
therefore certifies both notions.  This remains a local certificate: neither
notion implies global minimum weight.

\begin{proposition}[Certified pruning specification]
\label{prop:certified-pruning}
A pruning output is certified only together with its rechecked trace; the
trace certifies each deletion and endpoint productive-value preservation.
\end{proposition}

\begin{example}[Why deletion certificates must be rechecked]
\label{ex:t3-deletion-classes}
If two strategies are the only carriers of a module, each may be redundant
initially, but after deleting one the survivor is essential.  Applying both
stale certificates would remove the module from closure.
\end{example}

\subsection{Local versus global optimality}

\emph{Economic question.}  Does the absence of another certified
one-strategy deletion imply that the retained library has minimum burden?

A library is globally minimum weight precisely when it belongs to
\(\mathcal M_{\mathrm{safe}}(L)\) in
\eqref{eq:global-safe-optimizer}; it must be compared with every safe
sublibrary, not only its one-deletion neighbors.

\begin{theorem}[Local/global boundary]
\label{thm:local-global-safe-compression}
With strictly positive active weights, every global minimum is one-deletion
irreducible.  Every rechecked deletion trace is safe, but a one-deletion-
irreducible endpoint need not be globally minimum weight.  Even deleting the
unique heaviest currently safe strategy can lead to a strictly heavier
endpoint than another safe representation.
\end{theorem}

The theorem proves only the forward local implication: a global
minimum-weight safe representation cannot contain a safely deletable
positive-weight strategy.  Its converse is false, and the following exact
\((2,2,3)\) construction shows that even a uniquely determined
heaviest-safe-first step can terminate above the global minimum.

\begin{example}[Worked example: certified greedy endpoint versus global minimum]
\label{ex:worked-safe-greedy-global}
Take one belief, identity closure, and three zero-profile active strategies.
Strategies \(s_1\) and \(s_2\) carry \(\{m_1\}\) and \(\{m_2\}\) with weights
\(2\) and \(2\); the bundle \(s_{12}\) carries
\(\{m_1,m_2\}\) with weight \(3\).  The inactive strategy has zero profile,
no modules, and zero weight.  The current library is
\(L=\{s_0,s_1,s_2,s_{12}\}\).

\begin{center}
\small
% Generated by julia/scripts/generate_main_text_tables.jl.
% Exact source: experiments/results/resource_optimization_fixtures/02_cx_opt_prune_weight_01.json.
\begin{tabular}{@{}lcccc@{}}
\toprule
Representation & Active strategies & Burden (units) & Frontier (payoff units) & Closure \\
\midrule
Current library & \(s_{1},s_{2},s_{12}\) & \(7\) & \(F\equiv0\) & \(\{m_{1},m_{2}\}\) \\
Heaviest-safe-first endpoint & \(s_{1},s_{2}\) & \(4\) & \(F\equiv0\) & \(\{m_{1},m_{2}\}\) \\
Minimum-weight safe library & \(s_{12}\) & \(3\) & \(F\equiv0\) & \(\{m_{1},m_{2}\}\) \\
\bottomrule
\end{tabular}

\end{center}

All displayed values are exact rationals; denominator-one values are rendered
as integers.

At the source, \(s_{12}\) is the unique heaviest safely deletable strategy.
Deleting it first saves \(3\) units and leaves the inclusion-wise irreducible
pair \(\{s_1,s_2\}\); neither singleton can then be removed.  Complete
enumeration of the five safe sublibraries shows that
\(\{s_0,s_{12}\}\) is the minimum-weight safe representation.  It saves
\(4\) units relative to the current library, one more than the greedy endpoint.
Both endpoints have exactly the source frontier and closure, so their
productive values agree; only resource burden differs.  Deleting the two
singletons before the bundle reaches the global solution, making the order
dependence explicit.
\end{example}

The registered counterexample establishes the clean separation
\[
 \text{certified one-deletion irreducibility}
 \;\not\Longrightarrow\;
 \text{global minimum-weight compression}.
\]
It does not provide a greedy approximation ratio or imply deletion-order
invariance.  Appendix~B gives the exhaustive five-library comparison and the
general minimum-attainment and global-to-local arguments.  Those arguments
use finiteness for attainment and positive active weights for the
global-to-local implication; they use neither detectability nor probabilistic
independence.
